\sectionheading{5.22 ড্রাই রান ও ট্রেস টেবিল (Dry Run and Trace Table)}

Program বা Algorithm তৈরি করার পর Computer-এ Execute করার আগেই
কাগজে-কলমে ধাপে ধাপে পরীক্ষা করলে অনেক Logical Error সহজে
শনাক্ত করা যায়।

Program বাস্তবে Execute না করে হাতে ধাপে ধাপে পরীক্ষা করার
এই পদ্ধতিকে \textbf{Dry Run} বলা হয়।

Dry Run-এর সময় বিভিন্ন Variable ও Condition-এর পরিবর্তিত মান
সুশৃঙ্খলভাবে Table-এ লিখে রাখলে Program-এর কার্যপ্রবাহ আরও
সহজে বোঝা যায়। এই Table-ই হলো \textbf{Trace Table}।

%================================================
% Dry Run Definition
%================================================
\booksubheading{Dry Run কাকে বলে?}

\begin{definitionbox}

\textbf{\color{darkTeal}সংজ্ঞা}

কোনো Program বা Algorithm Computer-এ বাস্তবে Execute না করে
\textbf{নির্দিষ্ট Input ব্যবহার করে কাগজে-কলমে প্রতিটি ধাপ অনুসরণ
করে Result যাচাই করার প্রক্রিয়াকে Dry Run} বলে।

\end{definitionbox}

\textbf{সহজভাবে মনে রাখবে—}

\begin{center}

\textbf{
Program না চালিয়ে
\(\longrightarrow\)
হাতে ধাপ অনুসরণ
\(\longrightarrow\)
Result পরীক্ষা
=
Dry Run
}

\end{center}

%================================================
% Importance
%================================================
\booksubheading{Dry Run কেন করা হয়?}

Dry Run-এর মাধ্যমে—

\begin{compacttwocolpoints}

\item Program-এর Logic সঠিক কি না যাচাই করা যায়।

\item Variable-এর মান কীভাবে পরিবর্তিত হচ্ছে তা বোঝা যায়।

\item Condition কোন ক্ষেত্রে True বা False হচ্ছে তা দেখা যায়।

\item Loop কতবার চলছে তা অনুসরণ করা যায়।

\item Expected Output পাওয়া যাচ্ছে কি না যাচাই করা যায়।

\item Logical Error শনাক্ত করা সহজ হয়।

\end{compacttwocolpoints}

%================================================
% General Steps
%================================================
\booksubheading{Dry Run করার সাধারণ পদ্ধতি}

Dry Run করার সময় সাধারণত—

\begin{enumerate}

\item একটি Sample Input নেওয়া হয়।

\item Program-এর প্রথম Instruction থেকে শুরু করা হয়।

\item প্রতিটি Statement অনুযায়ী Variable-এর মান পরিবর্তন করা হয়।

\item Condition থাকলে True বা False নির্ণয় করা হয়।

\item Loop থাকলে প্রতিটি Iteration অনুসরণ করা হয়।

\item শেষে পাওয়া Output প্রত্যাশিত Result-এর সঙ্গে মিলিয়ে দেখা হয়।

\end{enumerate}

%================================================
% Trace Table
%================================================
\booksubheading{Trace Table}

Dry Run করার সময় Variable-এর মানগুলো এলোমেলোভাবে লিখলে
Program অনুসরণ করা কঠিন হতে পারে।

তাই প্রতিটি ধাপে Variable, Condition এবং Output-এর প্রয়োজনীয়
পরিবর্তিত মান Table-এ লিখে রাখা হয়।

%================================================
% Trace Table Definition
%================================================
\booksubheading{Trace Table কাকে বলে?}

\begin{definitionbox}

\textbf{\color{darkTeal}সংজ্ঞা}

Program বা Algorithm Dry Run করার সময়
\textbf{প্রতিটি ধাপ বা Iteration-এ গুরুত্বপূর্ণ Variable,
Condition ও Output-এর পরিবর্তিত মান সারণিতে লিপিবদ্ধ করার
পদ্ধতিকে Trace Table} বলে।

\end{definitionbox}

\begin{center}

\textbf{
Dry Run-এর ধাপগুলো Table-এ লিখে রাখা
=
Trace Table
}

\end{center}

%================================================
% Example 1
%================================================
\begin{examplebox}{উদাহরণ ১: সাধারণ হিসাবের Dry Run}

ধরা যাক—

\begin{codebox}

\textenglish{\ttfamily a = 5}

\textenglish{\ttfamily b = 3}

\textenglish{\ttfamily c = a + b}

\textenglish{\ttfamily c = c * 2}

\end{codebox}

Trace Table হবে—

\begin{tabularx}{\linewidth}{
|>{\RaggedRight\arraybackslash}Y|
>{\centering\arraybackslash}p{0.10\linewidth}|
>{\centering\arraybackslash}p{0.10\linewidth}|
>{\centering\arraybackslash}p{0.12\linewidth}|}
\hline

\rowcolor{softBlue}
\multicolumn{1}{|c|}{\textbf{ধাপ}}
&
\textbf{a}
&
\textbf{b}
&
\textbf{c}
\\
\hline

\textenglish{\ttfamily a = 5}
&
5
&
--
&
--
\\
\hline

\textenglish{\ttfamily b = 3}
&
5
&
3
&
--
\\
\hline

\textenglish{\ttfamily c = a + b}
&
5
&
3
&
8
\\
\hline

\textenglish{\ttfamily c = c * 2}
&
5
&
3
&
16
\\
\hline

\end{tabularx}

\vspace{1mm}

সুতরাং শেষ Result—

\begin{center}

\textbf{\textenglish{\ttfamily c = 16}}

\end{center}

Trace Table দেখে সহজেই বোঝা যায় কোন Statement-এর পরে
\textenglish{\ttfamily c}-এর মান কীভাবে পরিবর্তিত হয়েছে।

\end{examplebox}

%================================================
% Example 2 Loop
%================================================
\begin{examplebox}{উদাহরণ ২: Loop-এর Trace Table}

ধরা যাক 1 থেকে 3 পর্যন্ত সংখ্যার যোগফল নির্ণয়ের Logic—

\begin{codebox}

\textenglish{\ttfamily sum = 0}

\textenglish{\ttfamily for i = 1 to 3}

\hspace{8mm}\textenglish{\ttfamily sum = sum + i}

\end{codebox}

Trace Table—

\begin{tabularx}{\linewidth}{
|>{\centering\arraybackslash}p{0.20\linewidth}|
>{\centering\arraybackslash}p{0.15\linewidth}|
>{\centering\arraybackslash}Y|
>{\centering\arraybackslash}Y|}
\hline

\rowcolor{softBlue}
\textbf{Iteration}
&
\textbf{i}
&
\textbf{আগের sum}
&
\textbf{নতুন sum}
\\
\hline

1
&
1
&
0
&
1
\\
\hline

2
&
2
&
1
&
3
\\
\hline

3
&
3
&
3
&
6
\\
\hline

\end{tabularx}

\vspace{1mm}

অতএব—

\begin{center}

\textbf{Final sum = 6}

\end{center}

এই ধরনের Trace Table Loop-এর প্রতিটি Iteration বোঝার ক্ষেত্রে
বিশেষভাবে কার্যকর।

\end{examplebox}

%================================================
% Condition Example
%================================================
\booksubheading{Condition-এর ক্ষেত্রে Trace Table}

ধরা যাক—

\begin{center}

\textenglish{\ttfamily n = 8}
\end{center}

এবং পরীক্ষা করতে হবে—

\begin{center}

\textenglish{\ttfamily n \% 2 == 0}
\end{center}

তাহলে Trace Table—

\begin{tabularx}{\linewidth}{
|>{\centering\arraybackslash}p{0.16\linewidth}|
>{\centering\arraybackslash}Y|
>{\centering\arraybackslash}Y|}
\hline

\rowcolor{softBlue}
\textbf{n}
&
\textbf{Condition}
&
\textbf{Result}
\\
\hline

8
&
\textenglish{\ttfamily 8 \% 2 == 0}
&
True \(\rightarrow\) Even
\\
\hline

\end{tabularx}

\vspace{1mm}

অর্থাৎ Trace Table-এ শুধু Variable নয়,
প্রয়োজনে \textbf{Condition-এর Result}-ও লেখা যায়।

%================================================
% Columns
%================================================
\booksubheading{Trace Table-এ কী কী Column থাকবে?}

সব Program-এর জন্য একই Column প্রয়োজন হয় না।

Problem অনুযায়ী সাধারণত রাখা যেতে পারে—

\begin{compacttwocolpoints}

\item Step বা Iteration

\item Input Variable

\item Counter Variable

\item পরিবর্তিত Variable

\item Condition-এর Result

\item Output

\end{compacttwocolpoints}

\begin{rememberbox}

\textbf{\color{deepOrange}মনে রাখবে}

যেসব Variable বা Condition-এর মান পরিবর্তন অনুসরণ করা প্রয়োজন,
শুধু সেগুলো Trace Table-এ রাখলেই যথেষ্ট।

অপ্রয়োজনীয় Column যোগ করলে Table জটিল হয়ে যেতে পারে।

\end{rememberbox}

%================================================
% Dry Run vs Trace Table
%================================================
\booksubheading{Dry Run ও Trace Table-এর সম্পর্ক ও পার্থক্য}

Dry Run এবং Trace Table একই বিষয় নয়, তবে দুটি একে অপরের সঙ্গে
ঘনিষ্ঠভাবে সম্পর্কিত।

\begin{tabularx}{\linewidth}{|>{\bfseries}p{0.22\linewidth}|Y|Y|}
\hline

\rowcolor{softBlue}
\multicolumn{1}{|c|}{\textbf{বিষয়}}
&
\multicolumn{1}{c|}{\textbf{Dry Run}}
&
\multicolumn{1}{c|}{\textbf{Trace Table}}
\\
\hline

কী
&
Program হাতে ধাপে ধাপে পরীক্ষা করা
&
পরিবর্তিত মানগুলো Table-এ লেখা
\\
\hline

মূল উদ্দেশ্য
&
Logic ও Result যাচাই
&
মানের পরিবর্তন সহজে অনুসরণ
\\
\hline

রূপ
&
ধাপে ধাপে হাতে পরীক্ষা
&
সারণিভিত্তিক Record
\\
\hline

বিশেষ সুবিধা
&
Program-এর Flow বোঝা
&
Variable, Loop ও Condition পরিষ্কারভাবে দেখা
\\
\hline

\end{tabularx}

\begin{center}

\textbf{
Dry Run = পরীক্ষা করার প্রক্রিয়া
}

\vspace{0.5mm}

\textbf{
Trace Table = পরীক্ষার মান লিখে রাখার সারণি
}

\end{center}

%================================================
% Dry Run vs Debugging
%================================================
\booksubheading{Dry Run ও Debugging কি একই?}

না, দুটি একই বিষয় নয়।

\begin{tabularx}{\linewidth}{|>{\bfseries}p{0.24\linewidth}|Y|Y|}
\hline

\rowcolor{softBlue}
\multicolumn{1}{|c|}{\textbf{বিষয়}}
&
\multicolumn{1}{c|}{\textbf{Dry Run}}
&
\multicolumn{1}{c|}{\textbf{Debugging}}
\\
\hline

কাজ
&
Program না চালিয়ে Logic হাতে পরীক্ষা
&
Program-এর Error শনাক্ত ও সংশোধন
\\
\hline

Computer Execution
&
অবশ্যই প্রয়োজন হয় না
&
সাধারণত Program পরীক্ষা/চালানোর সঙ্গে সম্পর্কিত
\\
\hline

মূল লক্ষ্য
&
Logic ও Expected Output যাচাই
&
Error দূর করা
\\
\hline

\end{tabularx}

Dry Run Logical Error খুঁজে বের করতে এবং Debugging সহজ করতে
সহায়তা করতে পারে।

%================================================
% Quick Process
%================================================
\booksubheading{Dry Run-এর দ্রুত ধারা}

\begin{center}
\begin{adjustbox}{max width=0.98\linewidth}
\fcolorbox{darkBlue}{softBlue}{
    \parbox[c][9mm][c]{0.17\linewidth}{
        \centering\textbf{Sample\\Input}
    }
}
\(\longrightarrow\)
\fcolorbox{mainTeal}{softGreen}{
    \parbox[c][9mm][c]{0.18\linewidth}{
        \centering\textbf{Step\\Follow}
    }
}
\(\longrightarrow\)
\fcolorbox{darkBlue}{softBlue}{
    \parbox[c][9mm][c]{0.18\linewidth}{
        \centering\textbf{Value\\Change}
    }
}
\(\longrightarrow\)
\fcolorbox{mainTeal}{softGreen}{
    \parbox[c][9mm][c]{0.19\linewidth}{
        \centering\textbf{Output\\Check}
    }
}
\end{adjustbox}
\end{center}

%================================================
% Final Reminder
%================================================
\begin{rememberbox}

\textbf{\color{deepOrange}মনে রাখবে}

\begin{itemize}

\item Computer-এ Execute করার আগে হাতে Program বা Algorithm
পরীক্ষা করাকে \textbf{Dry Run} বলে।

\item Dry Run-এ Sample Input দিয়ে Statementগুলো ধারাবাহিকভাবে
অনুসরণ করা হয়।

\item Variable-এর মান কীভাবে পরিবর্তিত হচ্ছে তা দেখা
Dry Run-এর গুরুত্বপূর্ণ কাজ।

\item Dry Run-এর তথ্য সারণিতে লিখে রাখলে তাকে
\textbf{Trace Table} বলা হয়।

\item Trace Table বিশেষভাবে
\textbf{Loop, Variable ও Condition} বিশ্লেষণে কার্যকর।

\item Trace Table-এর Column Problem অনুযায়ী নির্ধারণ করতে হয়।

\item \textbf{Dry Run = Process}

\item \textbf{Trace Table = Record/Table}

\item Dry Run-এর প্রধান উদ্দেশ্য হলো
\textbf{Program-এর Logic ও Expected Output যাচাই করা}।

\item দ্রুত মনে রাখবে—

\begin{center}

\textbf{
Input
\(\rightarrow\)
Step Follow
\(\rightarrow\)
Value Track
\(\rightarrow\)
Output Verify
}

\end{center}

\end{itemize}

\end{rememberbox}

\small

%================================================
% 5.23 Algorithm
%================================================
